\chapterbackmatter{Supplementary Exercises}

\begin{enumerate}
\SupplementaryExercise{1}{M{\o}ller Operators and Scattering States and \(S\)-Matrix as an In--Out Overlap}
  {(Inspired by D. Harlow, \emph{Quantum Field Theory I}, 2024, Section 8.3 | D. Harlow, \emph{Quantum Field Theory I}, 2024, Sections 8.3--8.4.)}
  {harlow-qft1-2024-curated-parent-1}
\begin{enumerate}
\item
Define
\(\Omega_\pm=\lim_{t\to\mp\infty}e^{iHt}e^{-iH_0t}\), whenever the strong
limits exist.  Show that \(H\Omega_\pm=\Omega_\pm H_0\), identify
\(\InKet{\alpha}\) and \(\OutKet{\alpha}\), and write the scattering
operator in terms of \(\Omega_\pm\).
\item
Starting from translation covariance, prove that
\(\OutBra{\beta}\InKet{\alpha}\) can be nonzero only when the total
four-momenta agree.  Separate the identity contribution and explain why
the remaining connected matrix element contains a four-momentum delta
function.
\end{enumerate}

\SupplementaryExercise{2}{Completeness on the Scattering Subspace}
  {(Adapted from D. Harlow, \emph{Quantum Field Theory I}, 2024, Section 8.6, Problem 3.)}
  {harlow-qft1-2024-curated-parent-2}
For free identical scalar states normalized by
\(\braket{\mathbf q}{\mathbf p}=\delta^3(\mathbf q-\mathbf p)\), verify
the two-particle resolution of the identity, including its \(1/2!\)
factor.  State how asymptotic completeness transports this identity to
the in and out scattering bases.

\SupplementaryExercise{3}{Invariant Flux, Two-Body Phase Space, and Decay Widths}
  {(Inspired by D. Harlow, \emph{Quantum Field Theory I}, 2024, Sections 8.4 and 8.6, Problem 6 | D. Harlow, \emph{Quantum Field Theory I}, 2024, Section 8.4.)}
  {harlow-qft1-2024-curated-parent-3}
\begin{enumerate}
\item
For incident four-momenta \(p_1,p_2\), show that
\[
 u=\frac{\sqrt{(p_1\cdot p_2)^2-m_1^2m_2^2}}{E_1E_2}
\]
in the mostly-plus metric.  Verify that it becomes the ordinary speed of
particle \(1\) when particle \(2\) is at rest and becomes
\(|\mathbf v_1-\mathbf v_2|\) for collinear velocities.
\item
Starting with the four-dimensional delta function in
Eq.~\eqref{eq:3.4.15}, integrate the two-body final-state phase space in
the center-of-mass frame and rederive Eq.~\eqref{eq:3.4.28}, including
the unequal-mass expression for \(k'\).
\item
A spinless particle of mass \(M\) at rest decays to two nonidentical
spinless particles of masses \(m_1,m_2\), with angle-independent
Weinberg-normalized matrix element \(M_{fi}\).  Derive the total width and
state the change for two identical daughters.
\end{enumerate}

\SupplementaryExercise{4}{The Optical Theorem and Forward Scattering}
  {(Inspired by D. Harlow, \emph{Quantum Field Theory I}, 2024, Section 8.5.)}
  {harlow-qft1-2024-curated-parent-4}
\begin{enumerate}
\item
Write \(S=1-i\mathcal T\).  Derive the off-diagonal unitarity relation for
\(\mathcal T\), then set the external states equal and obtain the optical
theorem.  Explain why its right-hand side is nonnegative.
\item
Combine the optical theorem with Eq.~\eqref{eq:3.4.15} to express the
total cross section for a two-particle in-state in terms of the imaginary
part of its forward invariant amplitude.  Keep the book's state
normalization explicit.
\end{enumerate}

\SupplementaryExercise{5}{Exact Scattering from a Delta Potential}
  {(Adapted from D. Harlow, \emph{Quantum Field Theory I}, 2024, Section 8.6, Problem 4.)}
  {harlow-qft1-2024-curated-parent-5}
For \(H=p^2/(2m)+V_0\delta(x)\), \(V_0>0\), construct stationary
left- and right-incident in-states.  Determine the reflection and
transmission amplitudes and the \(2\times2\) \(S\)-matrix, and verify its
unitarity.

\SupplementaryExercise{6}{Disconnected Spectators}
  {(Inspired by D. Harlow, \emph{Quantum Field Theory I}, 2024, Section 8.4.)}
  {harlow-qft1-2024-curated-parent-6}
Suppose one outgoing particle has exactly the same species and momentum as
one incoming particle and does not interact.  Show how a spectator delta
function factors from the corresponding \(S\)-matrix element.  Explain
why connected amplitudes remove these extra singular distributions.

\SupplementaryExercise{7}{Mandelstam Variables in Massless Kinematics}
  {(Adapted from B. Mistlberger and K. Zhou, Stanford PHYS-330, Problem Set 5, 2022, Problem 1.)}
  {zhou-stanford-phys330-ps5-2022-curated-parent-1}
\begin{enumerate}
\item
In the mostly-plus metric define
\[
 s=-(p_1+p_2)^2,\quad t=-(p_1-p_3)^2,\quad
 u=-(p_1-p_4)^2
\]
for \(1+2\to3+4\).  Prove
\(s+t+u=m_1^2+m_2^2+m_3^2+m_4^2\).
\item
For massless \(2\to2\) scattering in the center-of-mass frame, express
\(s,t,u\) in terms of \(E_{\rm cm}\) and the scattering angle.  Determine
the physical ranges of \(t\) and \(u\).
\end{enumerate}

\SupplementaryExercise{8}{Yukawa Decay, Exchange, and Heavy-Field Matching}
  {(Adapted from B. Mistlberger and K. Zhou, Stanford PHYS-330, Problem Set 5, 2022, Problem 2.)}
  {zhou-stanford-phys330-ps5-2022-curated-parent-2}
\begin{enumerate}
\item
Let a real scalar \(\phi\) of mass \(M\) couple to a complex scalar
\(\psi\) of mass \(m\) through
\(\mathcal L_{\rm int}=-g\psi^*\psi\phi\).  For \(M>2m\), calculate
\(\Gamma(\phi\to\psi\psi^*)\) using covariantly normalized states.
\item
Find the leading invariant amplitude for
\(\psi(p_1)\psi^*(p_2)\to\psi(p_3)\psi^*(p_4)\).  Identify the two
channels and give the low-energy limit when every invariant is small
compared with \(M^2\).
\item
Use the classical equation of motion for the heavy field in the preceding
part to derive the leading local four-\(\psi\) interaction.  Check
that its tree amplitude agrees with the low-energy expansion of heavy
exchange, including the combinatorial factor.
\end{enumerate}

\SupplementaryExercise{9}{The Linear \(O(N)\) Sigma Model}
  {(Adapted from B. Mistlberger and K. Zhou, Stanford PHYS-330, Problem Set 5, 2022, Problem 4.)}
  {zhou-stanford-phys330-ps5-2022-curated-parent-3}
\begin{enumerate}
\item Quantize \(N\) real free scalar fields of equal mass and show that
their two-point function is diagonal in the \(O(N)\) indices.
\item Add
\(\mathcal L_{\rm int}=-(\lambda/4)(\Phi_i\Phi_i)^2\), and obtain the
four-scalar rule by differentiating with respect to labeled fields.
\item Calculate the center-of-mass total cross sections for
\(\Phi_1\Phi_2\to\Phi_1\Phi_2\),
\(\Phi_1\Phi_1\to\Phi_2\Phi_2\), and
\(\Phi_1\Phi_1\to\Phi_1\Phi_1\), including final-state symmetry factors.
\item Take \(m^2=-\mu^2<0\).  Find the vacuum manifold, choose a vacuum
in the \(N\)-direction, and introduce one radial field \(\sigma\) and
\(N-1\) fields \(\pi_a\).
\item Expand about that vacuum.  Determine both masses, the constant
term, and every cubic and quartic coupling.
\item List all propagators and vertices in the broken parameterization,
with their flavor tensors and numerical factors.
\item Find the total width for \(\sigma\) to decay into every allowed
two-pion state.  Evaluate the lifetime for \(N=3\), \(\lambda=0.1\),
and \(\mu=10\,\mathrm{GeV}\).
\item Sum the contact and three \(\sigma\)-exchange graphs for
\(\pi_i\pi_j\to\pi_k\pi_l\).
\item Show that the pion amplitude vanishes when one external pion is
taken soft on shell.
\item Add \(+a\Phi_N\) to the Lagrangian.  Determine the shifted-vacuum
condition and the pion mass through first order in \(a\).
\end{enumerate}

\SupplementaryExercise{10}{A Connected \(3\to3\) Tree Amplitude}
  {(Adapted from D. Tong, \emph{Quantum Field Theory: Example Sheet 4}, 2007, Problem 1.)}
  {tong-qft-sheet4-2007-curated-parent-1}
In real \(\lambda\phi^4/4!\) theory, use the Dyson series to identify the
connected \(3\to3\) contribution at order \(\lambda^2\).  For one chosen
partition of the external momenta, derive the internal propagator and
count the distinct tree topologies.

\SupplementaryExercise{11}{Exponentiation of Vacuum Bubbles}
  {(Adapted from D. Tong, \emph{Quantum Field Theory: Example Sheet 4}, 2007, Problem 2.)}
  {tong-qft-sheet4-2007-curated-parent-2}
Through order \(\lambda^2\) in real \(\phi^4\) theory, use Wick's theorem
to show that disconnected vacuum bubbles assemble into the first terms
of the exponential of connected vacuum bubbles.  Explain why normalized
scattering amplitudes contain no vacuum-bubble factor.

\SupplementaryExercise{12}{An \(O(3)\)-Symmetric Scattering Theory}
  {(Adapted from D. Tong, \emph{Quantum Field Theory: Example Sheet 4}, 2007, Problem 3.)}
  {tong-qft-sheet4-2007-curated-parent-3}
\begin{enumerate}
\item Starting from three degenerate free real scalars, derive their
propagator and explain why it is proportional to \(\delta_{ij}\).
\item For
\(\mathcal L_{\rm int}=-(\lambda/8)(\phi_i\phi_i)^2\), give all free and
interacting Feynman rules, including the invariant rank-four tensor.
\item Evaluate \(\phi_i\phi_j\to\phi_k\phi_l\) at tree level.
\end{enumerate}

\SupplementaryExercise{13}{Fermion and Antifermion Yukawa Scattering}
  {(Adapted from D. Tong, \emph{Quantum Field Theory: Example Sheet 4}, 2007, Problem 4.)}
  {tong-qft-sheet4-2007-curated-parent-4}
For
\[
 \mathcal L=\frac12(\partial\phi)^2-\frac12\mu^2\phi^2
 +\bar\psi(i\slashed\partial-m)\psi-\lambda\phi\bar\psi\psi ,
\]
\begin{enumerate}
\item compute \(\psi\psi\to\psi\psi\), retaining both scalar-exchange
graphs and deriving their relative sign from the external Fock states;
\item compute \(\psi\bar\psi\to\psi\bar\psi\), including elastic exchange
and annihilation.  Display the spinor order and Mandelstam channel of
each term.
\end{enumerate}

\SupplementaryExercise{14}{Three Routes to a Half-Line Fourier Transform}
  {(Adapted from A. Guth, MIT 8.323, Problem Set 4, 2008, Problem 2.)}
  {mit-823-s08-ps4-curated-parent-1}
\begin{enumerate}
\item Regard \(g(t)=\theta(t)e^{-i\omega_0t}\) as a tempered
distribution.  Pair its Fourier transform with
\[
 \varphi_\sigma(\omega)=\frac{1}{\sqrt{2\pi}\sigma}
 e^{-(\omega-\omega_1)^2/(2\sigma^2)}
\]
by transforming the test function first, and evaluate the half-line
integral.
\item Replace \(g(t)\) by \(g_\epsilon(t)=g(t)e^{-\epsilon t}\).
Perform the Gaussian pairing before taking \(\epsilon\to0^+\).
\item Instead truncate the time integral at \(T\).  Pair its ordinary
Fourier transform with the Gaussian and take \(T\to\infty\).  Prove
that all three routes give the same principal-value-plus-delta
distribution.
\end{enumerate}

\SupplementaryExercise{15}{Partial-Wave Unitarity and Cross-Section Bounds}
  {(Inspired by D. Harlow, \emph{Quantum Field Theory I}, 2024, Section 8.5 | D. Harlow, \emph{Quantum Field Theory I}, 2024, Sections 8.4--8.5.)}
  {harlow-qft1-2024-curated-parent-7}
\begin{enumerate}
\item
For spinless elastic scattering with
\[
 f(\theta)=\sum_{\ell=0}^\infty(2\ell+1)f_\ell P_\ell(\cos\theta),
\]
use \(S^\dagger S=1\) to prove
\(\operatorname{Im}f_\ell=k|f_\ell|^2\).  Parameterize the general
elastic solution by a real phase shift.
\item
Using the phase-shift form of the preceding part, derive
\[
 \sigma_{\rm tot}=\frac{4\pi}{k^2}
 \sum_{\ell=0}^\infty(2\ell+1)\sin^2\delta_\ell
\]
and the maximum elastic cross section carried by one partial wave.
\end{enumerate}

\SupplementaryExercise{16}{The \(K\)-Matrix and Inelastic Partial Waves}
  {(Inspired by D. Harlow, \emph{Quantum Field Theory I}, 2024, Section 8.5.)}
  {harlow-qft1-2024-curated-parent-8}
\begin{enumerate}
\item
For a single elastic partial wave, combine the Cayley transform
\(S=(1-iK)(1+iK)^{-1}\) with \(S=e^{2i\delta}\).  Find \(K\) in terms of
\(\delta\), locate the phase-shift resonance, and explain why a real
\(K\) preserves unitarity.
\item
When other channels are open, write an elastic diagonal element as
\(S_\ell=\eta_\ell e^{2i\delta_\ell}\).  Prove
\(0\leq\eta_\ell\leq1\) and derive the elastic, reaction, and total
partial-wave cross sections.
\end{enumerate}

\SupplementaryExercise{17}{Breit--Wigner Phase Motion}
  {(Inspired by B. Mistlberger and K. Zhou, Stanford PHYS-330, Problem Set 5, 2022, Problems 1--2.)}
  {zhou-stanford-phys330-ps5-2022-curated-parent-4}
Take
\[
 S_\ell(E)=\frac{E-E_R-i\Gamma/2}{E-E_R+i\Gamma/2}.
\]
Show that it is unitary on the real axis, determine a continuous
phase shift, and derive the Lorentzian form and narrow-width limit of
\(\sin^2\delta_\ell\).
\end{enumerate}
